\documentclass[10pt,a4paper]{amsart}
\usepackage[margin=19mm]{geometry}
\usepackage{amsmath,amssymb,mathtools}
\usepackage[scr=rsfs,cal=euler]{mathalfa}
\usepackage{xfrac}
\usepackage[unicode,hidelinks,bookmarks=false]{hyperref}
\newcommand{\ie}{i.e.\ }
\newcommand{\cf}{cf.\ }
\newcommand{\resp}{respectively\ }
\newcommand{\Subsection}{\subsection}
\providecommand{\nobreakdash}{\nobreak\hbox{-}}
\setlength{\emergencystretch}{2em}
\multlinegap0pt
\usepackage{amssymb}
\usepackage{bm}
\usepackage[shortlabels]{enumitem}
\newtheorem{proposition}{Proposition}
\renewcommand{\vec}[1]{\mathbf{#1}}
\title{Structure-aware divergences for comparing probability distributions}
\author{Rohit Sahasrabuddhe, Renaud Lambiotte}
\date{Source: arXiv:2603.22237v1 (CC BY 4.0)}
\begin{document}
\maketitle

\noindent\emph{Source and licence.} Structure-aware divergences for comparing probability distributions. Version arXiv:2603.22237v1 (CC BY 4.0), distributed by arXiv under the Creative Commons Attribution 4.0 International licence, http://creativecommons.org/licenses/by/4.0/, as recorded in the arXiv licence metadata for this exact version and shown on the abstract page https://arxiv.org/abs/2603.22237v1. Full licence notice: \texttt{licenses/arxiv-2603.22237-CC-BY-4.0.txt}.\par\smallskip

\noindent\textit{Editorial scope.} Counted unit: the article's proposition prop:1-d (source0.tex lines 520--527). The article's complete original proof of it is delivered with it (source0.tex lines 528--546, one \texttt{proof} environment of 19 lines). No mathematics is written, repaired or re-derived: the statement and the proof are the article's own lines, in the article's own notation, and no retained line is substituted or re-flowed.  No further source-local context is needed: every cross-reference of every retained line resolves inside the two delivered spans.  Nothing was omitted from any retained span, so the package carries no omission record.  No retained line contains a citation, so the wrapper carries no bibliography and no bibliography entry.  No retained line contains a figure environment, an image-inclusion command or drawing code, so no image is delivered and the package declares no asset dependency.  Editorial formatting only, in addition: an AMS \texttt{amsart} preamble that restates the article's own declarations and macros used by the retained lines, namely the environments \texttt{proposition} on separate counters, together with the standard packages those lines need.  The retained environments print in this document as Proposition 1 in order of appearance, whereas the article prints the same items with the numbering of its own full document.  The complete native source of the article as distributed in the arXiv e-print for this exact version is bundled unchanged as \texttt{sources/197008/source0.tex}.
\begin{proposition}
    Consider a finite metric space of strictly negative type with distance matrix $\vec{D}$ and the similarity matrix $\vec{Z}$ with entries 
    \begin{equation*}
        Z_{ij}=1-\frac{D_{ij}}{\max_{i,j}D_{ij}}.
    \end{equation*}
    $\vec{Z} \succ 0$ if $\vec{1}$ is an eigenvector of $\vec{D}$.
\label{prop:1-d}
\end{proposition}

\begin{proof}
In the following, we will use $\vec{D}$ to refer to the rescaled distance matrix. Rescaling does not change the sign of eigenvalues, so $\vec{D}$ remains of negative type. Since $\vec{D}$ is of strictly negative type,
\begin{equation}
    \vec{v}^\top \vec{Dv} < 0 \; \forall \; \vec{v} \in \left\{ \vec{v} \in \mathbb{R}^n \mid \vec{1}^\top\vec{v}=0 \right\} \setminus \left\{ \vec{0}\right\}.
\end{equation}

For $\vec{Z}=\vec{11}^\top - \vec{D}$, we aim to show that
\begin{equation}
    \vec{x}^\top \vec{Zx} > 0 \; \forall\; \vec{x} \in \mathbb{R}^n \setminus \left\{\vec{0}\right\}.
\end{equation}

We begin by noting that $\left\{ \vec{v} \in \mathbb{R}^n \mid \vec{1}^\top\vec{v}=0 \right\} = \textnormal{span}\left(\vec{1}\right)^\perp$. Thus, we can write any $\vec{x}\in\mathbb{R}^n$ as $\vec{x}=\vec{v} + a\vec{1}$ for $\vec{v}\in\textnormal{span}\left(\vec{1}\right)^\perp$ and $a\in\mathbb{R}$. Substituting this in the quadratic form above, we get the following.
\begin{align*}
    \vec{x}^\top \vec{Zx} &= \left( \vec{v} + a\vec{1}\right)^\top \left( \vec{11}^\top - \vec{D} \right) \left( \vec{v} + a\vec{1} \right)\\
    &= \vec{v}^\top \vec{Z} \vec{v} + 2a \vec{v}^\top\vec{Z} \vec{1} + a^2\vec{1}^\top \vec{Z} \vec{1}
\end{align*}

Since $\vec{v}^\top \vec{11}^\top \vec{v}=0$, the first term equals $-\vec{v}^\top \vec{D} \vec{v} > 0$ since $\vec{D}$ is of strictly negative type. $\vec{1}^\top \vec{Z} \vec{1} \geq n$ since $Z_{ii}=1$ and $Z_{ij}\geq 0$. Therefore, $a^2\vec{1}^\top \vec{Z} \vec{1} \geq na^2 \geq 0$. Thus, when the cross-term $2a \vec{v}^\top\vec{Z} \vec{1}$ is non-negative, $\vec{x}^\top \vec{Zx} > 0$ as desired. When $\vec{1}$ is an eigenvalue of $\vec{Z}$, we can substitute $\vec{Z1}=\lambda \vec{1}$ in the cross-term to get $2a\lambda \vec{v}^\top\vec{1} = 0$. This concludes the proof.
\end{proof}

\end{document}
